% ═══════════════════════════════════════════════════════════════
%  chapters/ch01-intro.tex
%  Вводная глава: о проекте, о книге, о читателе
% ═══════════════════════════════════════════════════════════════

\chapter{Зачем два STM32?}
\label{ch:intro}

\epigraph{<<Осциллограф — это глаза инженера.
Генератор — это его голос.>>}{}

\section{О чём эта книга}
\label{sec:intro-about}

Эта книга — о том, как построить \textbf{два} устройства
на \textbf{микроконтроллерах STM32}, которые работают
\emph{вместе}, и как написать для них \textbf{программу
управления} на \textbf{GTK4}. Речь идёт о
\textbf{цифровом осциллографе} и \textbf{генераторе сигналов}~---
двух приборах, без которых не обходится ни одна
радиолаборатория.

\begin{ideabox}
\textbf{Это вторая книга серии.} Первая~--- <<USB, Linux и GTK>>~---
рассказывала про \emph{один} микроконтроллер и \emph{один}
последовательный порт. Здесь~--- \emph{два} микроконтроллера,
\emph{два} USB-устройства, \emph{бинарный} протокол,
\emph{многопоточность} и GTK4. Мы предполагаем, что вы
\emph{прочитали} первую книгу, но если нет~--- ничего страшного.
Мы будем \emph{кратко} напоминать ключевые моменты.
\end{ideabox}

Это не учебник <<в теории>>. Это книга о \textbf{настоящем}
проекте~--- со всеми его граблями, странностями и
инженерными компромиссами. Мы будем писать код, который
\emph{работает}, и объяснять, \emph{почему} он работает
именно так.

Вот что мы сделаем:

\begin{itemize}
    \item \textbf{Осциллограф}~--- на плате \textbf{NUCLEO-F401RE}.
          Он захватывает сигнал через ADC, обрабатывает его
          через DMA и передаёт на компьютер через USB.
    \item \textbf{Генератор}~--- на плате \textbf{STM32F103C8T6}
          (Blue Pill). Он формирует сигнал через DAC или PWM
          с использованием DMA и таблиц форм.
    \item \textbf{Программа управления}~--- на \textbf{GTK4}.
          Она принимает данные с осциллографа, управляет
          генератором и рисует осциллограмму в реальном времени.
\end{itemize}

\begin{keybox}
\textbf{Всё вместе}~--- это \emph{миниатюрная радиолаборатория}
на вашем столе. Осциллограф, генератор и программа
управления. Стоимость железа~--- \emph{несколько тысяч}
рублей. Стоимость знаний, которые вы получите~--- \emph{бесценна}.
\end{keybox}

\section{Почему два микроконтроллера, а не один}
\label{sec:intro-two-mcu}

Первый вопрос, который возникает: зачем \emph{два} STM32?
Почему не \emph{один}, который делает всё? Ответ~---
\textbf{ресурсы} и \textbf{разделение ответственности}.

\subsection{Ресурсы}

Каждая задача~--- осциллограф и генератор~---
сама по себе \emph{требовательна} к ресурсам микроконтроллера.
Осциллографу нужны:

\begin{itemize}
    \item \textbf{Быстрый ADC}~--- чтобы захватывать сигнал
          с частотой до \SI{500}{\kilo\hertz}.
    \item \textbf{DMA}~--- чтобы перекладывать данные из ADC
          в память \emph{без участия CPU}.
    \item \textbf{Двойной буфер} (ping-pong)~--- чтобы
          \emph{непрерывно} захватывать сигнал, не теряя данные.
    \item \textbf{USB CDC}~--- чтобы передавать данные
          на компьютер.
\end{itemize}

Генератору нужны:

\begin{itemize}
    \item \textbf{Точный таймер}~--- чтобы формировать сигнал
          с заданной частотой (от \SI{1}{\hertz} до \SI{100}{\kilo\hertz}).
    \item \textbf{DAC} или \textbf{PWM}~--- чтобы выдавать
          аналоговый сигнал.
    \item \textbf{DMA circular}~--- чтобы \emph{непрерывно}
          воспроизводить таблицу формы.
    \item \textbf{USB CDC}~--- чтобы принимать команды
          с компьютера.
\end{itemize}

\begin{warnbox}
\textbf{На одном STM32F103 это \emph{можно} сделать~--- но
с трудом.} У него \emph{один} ADC, \emph{один} DAC,
\emph{один} USB. Придётся \emph{жонглировать} прерываниями
и таймерами. Мы получим \emph{нестабильный} прибор,
в котором осциллограф <<заикается>>, когда генератор меняет
частоту.
\end{warnbox}

\subsection{Разделение ответственности}

\textbf{Разделение ответственности}~--- это принцип, знакомый
любому программисту. \emph{Каждый модуль} делает \emph{одну}
вещь и делает её хорошо. То же самое~--- в железе.

\begin{ideabox}
\textbf{Аналогия: оркестр.} В симфоническом оркестре \emph{не}
один музыкант играет \emph{все} партии. Есть скрипка, есть
виолончель, есть флейта. Каждый инструмент~--- \emph{своя}
задача. То же с STM32: один захватывает, другой формирует.
\end{ideabox}

\subsection{Бонус: отладка проще}

Когда у вас \emph{два} независимых устройства~--- их
\emph{проще} отлаживать. Если осциллограф <<глючит>>~---
вы \emph{точно} знаете, что проблема в \emph{его} прошивке,
а не в том, что генератор <<отъел>> половину таймеров.

\section{Что такое осциллограф}
\label{sec:intro-oscilloscope}

\textbf{Осциллограф}~--- это прибор для \emph{визуализации}
электрических сигналов. Он показывает, как \emph{меняется}
напряжение во времени.

\begin{keybox}
\textbf{Осциллограф~--- это микроскоп для электричества.}
Мультиметр показывает \emph{одно} число~--- например,
<<\SI{5}{\volt}>>. Осциллограф показывает \emph{график}~---
как напряжение меняется \emph{каждую микросекунду}.
\end{keybox}

\subsection{Зачем он нужен}

Представьте, что у вас есть схема, которая должна
выдавать \emph{прямоугольный} сигнал частотой
\SI{1}{\kilo\hertz}. Вы подключаете мультиметр~---
он показывает \SI{2.5}{\volt}. Всё в порядке?

\textbf{Нет!} Мультиметр показывает \emph{среднее}
напряжение. А что если сигнал на самом деле~---
\emph{синус}, а не \emph{прямоугольник}? Что если есть
\emph{шум}, \emph{дребезг}, \emph{искажения}?
Мультиметр \emph{не покажет} этого. Осциллограф~---
\textbf{покажет}.

\subsection{Как он работает}

\emph{Цифровой} осциллограф (DSO~--- Digital Storage
Oscilloscope) работает так:

\begin{enumerate}
    \item \textbf{Входной сигнал} через \emph{делитель}
          попадает на \textbf{ADC}.
    \item \textbf{ADC} \emph{семплирует} сигнал с заданной
          частотой (например, \SI{1}{\mega\hertz}).
    \item \textbf{DMA} перекладывает семплы в память.
    \item \textbf{Процессор} (или ПК) \emph{рисует} их
          на экране.
\end{enumerate}

\subsection{Что мы сделаем}

Наш осциллограф на \textbf{NUCLEO-F401RE}:

\begin{itemize}
    \item \textbf{ADC}~--- \SI{12}{\bit}, до \SI{2.4}{\mega\sample/\second}.
    \item \textbf{DMA}~--- с \emph{двойным} буфером (ping-pong).
    \item \textbf{Триггер}~--- по уровню и фронту.
    \item \textbf{USB CDC}~--- передача кадров на ПК.
    \item \textbf{До \SI{500}{\kilo\hertz}} частоты дискретизации.
\end{itemize}

\begin{notebox}
\textbf{Это не замена HANTEK'у.} Наш осциллограф~---
\emph{учебный}. Он работает на частотах до
\SI{500}{\kilo\hertz}, тогда как ваш HANTEK~--- до
\SI{150}{\mega\hertz}. Но \textbf{принципы}~--- те же.
Поняв, как работает наш, вы поймёте, как работает любой.
\end{notebox}

\section{Что такое генератор}
\label{sec:intro-generator}

\textbf{Генератор сигналов}~--- это прибор, который
\emph{создаёт} электрические сигналы заданной формы,
частоты и амплитуды.

\begin{keybox}
\textbf{Генератор~--- это голос инженера.} Осциллограф
\emph{слушает}, генератор \emph{говорит}. Вместе они
позволяют \emph{исследовать} любую схему~---
подать сигнал, посмотреть, что получилось.
\end{keybox}

\subsection{Зачем он нужен}

Когда вы \emph{проектируете} схему, вам нужно \emph{проверить},
как она работает. Для этого нужен \emph{известный} сигнал
на входе. Генератор даёт этот сигнал.

Пример: вы сделали \emph{фильтр низких частот}.
Подайте на вход \emph{синус} \SI{1}{\kilo\hertz}~---
должен пройти. Подайте \SI{100}{\kilo\hertz}~---
должен быть подавлен. Генератор~--- ваш \emph{инструмент
проверки}.

\subsection{Формы сигналов}

Стандартные формы:

\begin{itemize}
    \item \textbf{Синус}~--- \emph{гладкая} кривая, основная
          форма в природе.
    \item \textbf{Меандр}~--- \emph{прямоугольник},
          используется в цифровых схемах.
    \item \textbf{Пила}~--- \emph{линейно} растёт и
          \emph{резко} падает, используется в развёртках.
    \item \textbf{Треугольник}~--- \emph{линейно} растёт
          и \emph{линейно} падает.
    \item \textbf{Пользовательская}~--- \emph{любая} форма,
          загруженная пользователем.
\end{itemize}

\subsection{Что мы сделаем}

Наш генератор на \textbf{STM32F103C8T6}:

\begin{itemize}
    \item \textbf{DAC}~--- \SI{12}{\bit}, до \SI{1}{\mega\sample/\second}.
    \item \textbf{PWM}~--- альтернатива DAC для \emph{меандра}.
    \item \textbf{DMA circular}~--- \emph{непрерывное}
          воспроизведение таблицы.
    \item \textbf{Частота}~--- от \SI{1}{\hertz} до \SI{100}{\kilo\hertz}.
    \item \textbf{Формы}~--- синус, меандр, пила, треугольник,
          двухполупериодный, однополупериодный, пользовательская.
\end{itemize}

\section{Что умеет наша программа}
\label{sec:intro-capabilities}

Программа управления (ПК-приложение на GTK4) позволяет:

\begin{itemize}
    \item \textbf{Выбирать порты} осциллографа и генератора.
    \item \textbf{Запускать и останавливать} поток данных
          с осциллографа.
    \item \textbf{Настраивать генератор}: форму, частоту,
          амплитуду, смещение, скважность.
    \item \textbf{Рисовать осциллограмму} в реальном времени.
    \item \textbf{Сохранять} данные в файл.
\end{itemize}

\begin{ideabox}
\textbf{Почему GTK4, а не GTK3?} Первая книга использовала
GTK3. GTK4~--- \emph{современная} версия, с \emph{новым}
API отрисовки, \emph{лучшей} поддержкой Wayland и
\emph{меньшим} потреблением памяти. Если вы знаете GTK3,
GTK4 покажется \emph{похожим}, но с \emph{нюансами}.
Мы \emph{разберём} их в главе \ref{ch:gtk4-basics}.
\end{ideabox}

\section{Что понадобится: железо}
\label{sec:intro-hardware}

Для работы с книгой понадобится:

\subsection{Обязательно}

\begin{itemize}
    \item \textbf{NUCLEO-F401RE}~--- плата с STM32F401.
          \emph{Осциллограф}.
    \item \textbf{STM32F103C8T6} (Blue Pill)~--- плата
          с STM32F103. \emph{Генератор}.
    \item \textbf{ST-Link V2}~--- программатор для Blue Pill.
          \emph{Nucleo} имеет \emph{встроенный} ST-Link.
    \item \textbf{Компьютер} на Linux (Debian 13, LMDE 7
          или аналог).
\end{itemize}

\subsection{Настоятельно рекомендуется}

\begin{itemize}
    \item \textbf{HANTEK DSO2D15} (или любой DSO)~---
          \emph{эталонный} осциллограф.
    \item \textbf{FG-200 DDS} (или любой генератор)~---
          \emph{эталонный} генератор.
    \item \textbf{Макетная плата} и провода.
    \item \textbf{Резисторы} и конденсаторы (для RC-фильтра).
\end{itemize}

\subsection{Опционально}

\begin{itemize}
    \item \textbf{USBASP 2.0}~--- для AVR (если хотите
          перепрошить Arduino без bootloader).
    \item \textbf{FINRSI DSO-TC4}~--- карманный осциллограф.
\end{itemize}

\begin{ideabox}
\textbf{Почему нужны эталонные приборы?} Наш осциллограф и
генератор~--- \emph{учебные}. Чтобы понять, \emph{правильно}
ли они работают, нужно сравнить с \emph{эталоном}. Если
у вас нет HANTEK и FG-200~--- можно \emph{сравнить} наш
осциллограф с нашим генератором: подайте сигнал с генератора
на осциллограф и проверьте, совпадает ли форма.
\end{ideabox}

\section{Что понадобится: софт}
\label{sec:intro-software}

Мы работаем на \textbf{Debian 13} (или \textbf{LMDE 7}).
Все инструменты доступны через \texttt{apt}.

\subsection{Основные инструменты}

\begin{minted}[frame=none, bgcolor=gray!5]{bash}
sudo apt install gcc-arm-none-eabi binutils-arm-none-eabi \
                 gdb-multiarch stlink-tools openocd \
                 build-essential pkg-config git
\end{minted}

\subsection{Библиотеки для STM32}

\begin{itemize}
    \item \textbf{STM32CubeF4}~--- HAL для STM32F401 (Nucleo).
    \item \textbf{STM32CubeF1}~--- HAL для STM32F103 (Blue Pill).
\end{itemize}

\begin{minted}[frame=none, bgcolor=gray!5]{bash}
cd ~/opt
git clone https://github.com/STMicroelectronics/STM32CubeF4.git
git clone https://github.com/STMicroelectronics/STM32CubeF1.git
\end{minted}

\subsection{Библиотеки для ПК}

\begin{minted}[frame=none, bgcolor=gray!5]{bash}
sudo apt install libgtk-4-dev libglib2.0-dev libcairo2-dev
\end{minted}

\begin{warnbox}
\textbf{Пакет \texttt{gdb-arm-none-eabi} больше не существует.}
Начиная с Debian 10, он \emph{удалён}. Вместо него~---
\texttt{gdb-multiarch}. Чтобы команда \texttt{arm-none-eabi-gdb}
работала, создайте симлинк:
\begin{minted}{bash}
sudo ln -s /usr/bin/gdb-multiarch /usr/bin/arm-none-eabi-gdb
\end{minted}
Это \emph{реальная} грабля~--- вы столкнётесь с ней, если
попробуете установить \texttt{gdb-arm-none-eabi}.
\end{warnbox}

\section{Что вы узнаете}
\label{sec:intro-structure}

Книга состоит из \textbf{семи частей}:

\subsection{Часть I. Введение}

\begin{itemize}
    \item \textbf{Глава 1}~--- эта глава.
    \item \textbf{Глава 2}~--- \emph{прошивка STM32 из
          командной строки}. Инструменты, ST-Link,
          первая программа (мигание LED).
\end{itemize}

\subsection{Часть II. Архитектура и протокол}

\begin{itemize}
    \item \textbf{Глава 3}~--- \emph{архитектура} проекта.
          Три мира: осциллограф, генератор, ПК.
    \item \textbf{Глава 4}~--- \emph{протокол} обмена.
          Кадры, CRC16, seq, команды.
\end{itemize}

\subsection{Часть III. Осциллограф}

\begin{itemize}
    \item \textbf{Глава 5}~--- \emph{ADC + DMA}. Как
          захватить сигнал.
    \item \textbf{Глава 6}~--- \emph{триггер}. Как найти
          начало сигнала.
    \item \textbf{Глава 7}~--- \emph{кольцо кадров}.
          Как \emph{непрерывно} передавать данные.
\end{itemize}

\subsection{Часть IV. Генератор}

\begin{itemize}
    \item \textbf{Глава 8}~--- \emph{DAC, PWM, RC-фильтр}.
    \item \textbf{Глава 9}~--- \emph{таблицы форм}.
          DDS на STM32.
    \item \textbf{Глава 10}~--- \emph{DMA circular}.
          Частота, амплитуда, смещение.
\end{itemize}

\subsection{Часть V. ПК-приложение}

\begin{itemize}
    \item \textbf{Глава 11}~--- \emph{GTK4}. Что нового
          по сравнению с GTK3.
    \item \textbf{Глава 12}~--- \emph{архитектура ПК-приложения}.
          Потоки, мьютексы, очередь кадров.
    \item \textbf{Глава 13}~--- \emph{отрисовка осциллограммы}.
          Cairo для быстрых данных.
    \item \textbf{Глава 14}~--- \emph{интерфейс генератора}.
          Формы, частота, амплитуда.
\end{itemize}

\subsection{Часть VI. Завершение}

\begin{itemize}
    \item \textbf{Глава 15}~--- \emph{жизненный цикл}.
          Запуск, работа, завершение.
    \item \textbf{Глава 16}~--- \emph{отладка}. gdb,
          valgrind, OpenOCD.
    \item \textbf{Глава 17}~--- \emph{AppImage}. Как
          упаковать приложение.
\end{itemize}

\subsection{Приложения}

\begin{itemize}
    \item \textbf{A}~--- исходный код.
    \item \textbf{B}~--- сборка проекта.
    \item \textbf{C}~--- глоссарий.
\end{itemize}

\section{История проекта}
\label{sec:intro-history}

Эта книга родилась из реального опыта. Однажды студент
показал преподавателю прототип осциллографа на STM32.
Код работал, но был <<сырым>>: всё в одном файле,
без модулей, без объяснений. Преподаватель решил
разобраться и \emph{освоить} STM32 с нуля. Мы сели
и прошли весь путь вместе~--- от установки инструментов
до прошивки первой программы.

\emph{Грабли}, которые мы нашли, стали главами этой
книги. Мы \emph{специально} не причёсывали их, не убирали
<<неудобные>> моменты. Реальная инженерия~--- это когда
что-то не работает, и ты \emph{разбираешься}, почему.

\begin{ideabox}
\textbf{Эта книга~--- про \emph{путь}, а не про результат.}
Результат~--- работающий прибор~--- важен. Но \emph{путь}~---
важнее. Мы надеемся, что читатель пройдёт этот путь вместе
с нами и \emph{поймёт}, как устроена реальная разработка.
\end{ideabox}

\section{Соглашения}
\label{sec:intro-conventions}

\subsection{Типографика}

\begin{itemize}
    \item \texttt{monospace}~--- \textbf{код}, имена файлов,
          команды, идентификаторы.
    \item \textit{курсив}~--- \emph{термины} при первом
          упоминании, названия книг.
    \item \textbf{жирный}~--- \textbf{важные} мысли, ключевые
          понятия.
\end{itemize}

\subsection{Врезки}

Как и в первой книге, используем четыре вида врезок:

\begin{itemize}
    \item \faLightbulb~--- \textbf{Идея}. Общая мысль.
    \item \faKey~--- \textbf{Ключевой момент}. Важно запомнить.
    \item \faExclamationTriangle~--- \textbf{Внимание}.
          Грабля.
    \item \faInfoCircle~--- \textbf{Заметка}. Дополнительно.
\end{itemize}

\subsection{Структура глав}

Каждая глава построена одинаково:

\begin{enumerate}
    \item \textbf{Введение}~--- что разберём.
    \item \textbf{Архитектура}~--- как модуль вписывается
          в проект.
    \item \textbf{Разбор кода}~--- по функциям.
    \item \textbf{Диаграммы}~--- визуализация.
    \item \textbf{Уроки главы}~--- сводка.
    \item \textbf{Чек-лист грабель}~--- что может пойти не так.
    \item \textbf{Упражнения}~--- для закрепления.
    \item \textbf{Проверьте сами}~--- практические задания.
\end{enumerate}

\section{Исходный код}
\label{sec:intro-source}

Весь исходный код проекта \textbf{доступен} на GitHub:

\begin{center}
\url{https://github.com/danayer/dualboard-stm32-oscilloscope}
\end{center}

\textbf{Структура репозитория:}

\begin{verbatim}
dualboard-stm32-oscilloscope/
├── firmware/
│   ├── oscilloscope/       # Прошивка осциллографа
│   └── generator/          # Прошивка генератора
├── pc-app/                 # GTK4 приложение
│   ├── main.c
│   ├── Makefile
│   └── ...
├── docs/
│   ├── protocol.md         # Протокол обмена
│   └── ...
└── README.md
\end{verbatim}

\section{Проверьте сами}
\label{sec:intro-check}

Прежде чем читать дальше~--- \emph{проверьте}, что
у вас есть всё необходимое.

\textbf{1. Подключите Nucleo к компьютеру.}

\textbf{2. Выполните:}

\begin{minted}[frame=none, bgcolor=gray!5]{bash}
st-info --probe
\end{minted}

\textbf{Вы должны увидеть}:

\begin{minted}[frame=none, bgcolor=gray!5]{text}
Found 1 stlink programmers
  version:    V2J19S3
  serial:     066CFF484849785087031235
  flash:      524288 (pagesize: 16384)
  sram:       98304
  chipid:     0x433
  dev-type:   STM32F401xD_xE
\end{minted}

\begin{warnbox}
\textbf{Если видите два устройства}~--- отключите лишние.
Одновременная работа с двумя ST-Link~--- \emph{частая}
причина загадочных ошибок. Подробно~--- в главе~\ref{ch:flashing}.
\end{warnbox}

\textbf{3. Прочитайте Flash:}

\begin{minted}[frame=none, bgcolor=gray!5]{bash}
st-flash read /tmp/nucleo.bin 0x08000000 256
hexdump -C /tmp/nucleo.bin | head -20
\end{minted}

\textbf{Если видите мусор} (\texttt{[Trash Info]}, ASCII-текст)~---
\emph{не паникуйте}. Это норма для б/у платы. Мы разберём
в главе~\ref{ch:flashing}.

\begin{ideabox}
\textbf{Зачем это нужно?} Пока вы не убедились, что
ST-Link \emph{видит} плату, двигаться дальше нет смысла.
99\% проблем~--- именно здесь.
\end{ideabox}

\section{Что дальше}
\label{sec:intro-next}

В \textbf{следующей главе} мы разберём \textbf{прошивку STM32
из командной строки}: установим инструменты, настроим ST-Link,
разберёмся с типичными ошибками и прошьём \emph{первую}
программу~--- мигание LED.

\textbf{Если вы готовы}~--- переходите к главе~\ref{ch:flashing}.